\section{Hexadic--octadic incidence, area operator and sectoral information}
\label{rev4:ae:seccion}

The geometric recovery of the record acquires physical content when
composed with observables whose construction preserves the incidences and
state conditions used. The passage from a hexad to its marked octad
provides a precise case: enlarging the support transforms two
incidence spaces, determines their multiplicities and allows
that distinction to be preserved in an area operator and in the logarithm of a
reduced state. The following exposition brings together that construction from the second
article in the series~\cite{hmtII} and the marked Grassmannian
representation of the preceding section.

The starting point remains the enriched state produced by
APP--TRIT--TPK. The two APP operations preserve their complete
divisions, TRIT determines regime and orientation, and the composition
\(U_t=\operatorname{Upd}_t\circ\operatorname{Tra}_t\circ
\operatorname{Sel}_t\) transports and updates the information on position,
carry and memory. Compatible prolongation and nonadic holonomy
jointly preserve the arithmetic and incidence realizations of the
discrete structure of the continuum. The record \(K\), the regional
coordinates \(\pi_{\rm HMT},\varphi_{\rm HMT},e_{\rm HMT}\), the
coordinate \(\alpha_{\rm HMT}\) and the exceptional configuration arise
from that common construction. Here we use its results and the markings
that allow them to be prolonged towards boundary observables.

\subsection{Incidence lifting and difference of multiplicities}

Let \(\mathcal H\) be the Witt hexad system on the twelve positions
of the record. Fix the incidence chart \(\ell:[12]\to\mathscr D\),
where \(\mathscr D\) is a dodecad of the extended binary Golay code
\(G_{24}\), with parameters \([24,12,8]_2\). The chart transports the
hexads of \(\mathcal H\) to the residual design on \(\mathscr D\).
This realization datum preserves the marked configuration determined
in the exceptional development. The weight-eight supports of \(G_{24}\)
constitute the octads of the system \(S(5,8,24)\).

\begin{proposition}[Unique lift of a residual hexad]
\label{rev4:ae:elevacion}
The sets \(O\cap\mathscr D\), with \(O\) an octad and
\(|O\cap\mathscr D|=6\), form a system \(S(5,6,12)\).
Each hexad of that system is the intersection of \(\mathscr D\) with
a unique octad. In particular, the marked hexad \(H\subset\mathscr D\)
has a determined lift \(O_H\) and an inclusion \(H\subset O_H\).
\end{proposition}
\begin{proof}
The code weights are \(0,8,12,16,24\). If
\(j=|O\cap\mathscr D|\), the binary sum of the vectors with supports
\(O\) and \(\mathscr D\) has weight \(20-2j\); hence,
\(j\in\{2,4,6\}\). Every five-element subset \(T\subset\mathscr D\)
belongs to a unique octad. Since its intersection with \(\mathscr D\)
contains \(T\), that intersection has six elements. This yields
a unique residual hexad on each five-element subset. If two octads lifted
the same hexad, both would contain any of its five-element subsets and
would coincide by the Steiner property.
\end{proof}

In what follows, \(\binom H2\) denotes the set of two-element
subsets of \(H\). We preserve that second factor when enlarging the first:
\begin{equation}
 \mathcal F_6=H\times\binom H2,
 \qquad
 \mathcal F_8=O_H\times\binom H2,
 \qquad
 \jmath(h,P)=(h,P).
 \label{rev4:ae:incidencias}
\end{equation}
An element of these spaces is a point--pair incidence. The distinguished
point ranges first over the hexad and then over its octad, while the
pair continues to belong to the original hexad. This specification
fixes the operation responsible for the increase.

\begin{proposition}[Exact enlargement of the incidence space]
\label{rev4:ae:incremento}
The map \(\jmath\) is injective and its complement is
\((O_H\setminus H)\times\binom H2\). Consequently,
\begin{equation}
 |\mathcal F_6|=6\binom62=90,
 \qquad |\mathcal F_8|=8\binom62=120,
 \qquad |\mathcal F_8\setminus\jmath(\mathcal F_6)|=30.
 \label{rev4:ae:cardinales}
\end{equation}
\end{proposition}
\begin{proof}
The inclusion \(H\subset O_H\) preserves both coordinates and proves
injectivity. The disjoint partition \(O_H=H\sqcup(O_H\setminus H)\)
induces the partition of the Cartesian product. Since
\(|O_H\setminus H|=2\) and \(\binom62=15\), the complement has
\(2\cdot15=30\) elements.
\end{proof}

The two multiplicities are obtained before any probability
distribution on the boundary. The normalized difference
\((90-120)/(24\binom62)=-1/12\) preserves the sign of the oriented
comparison. Its interpretation depends on the incidence space in which
the operation has been performed.

\subsection{Collective transport and the Barbero--Immirzi parameter}

The boundary representation retains the labels of both sectors.
To make its support explicit, write
\(x_j=r_{2j-1}+3r_{2j}\), with \(r_k\in\{0,1,2\}\), in three
coordinates of \(\mathbb Z/9\mathbb Z\). The positional regrouping
\begin{equation}
 \beta_{729}(x_1,x_2,x_3)
 =\bigl(r_1+3r_2+9r_3,\ r_4+3r_5+9r_6\bigr)
 \label{rev4:ae:reagrupamiento}
\end{equation}
is a bijection onto \((\mathbb Z/27\mathbb Z)^2\). Uniqueness of
the ternary expansion proves that its inverse extracts the six digits and
restores the three initial pairs. Positional order is preserved
with its carry information; the respective additions remain
defined in their own groups.

In the block \(B_\partial=\{0,\ldots,8\}^2\) of the torus of order
twenty-seven, we distinguish the oriented links
\begin{align}
 E_{90}&=\{((-1,j),(0,j)),\ ((8,j),(9,j)):0\le j\le8\},\\
 E_{120}&=\{((i,-1),(i,0)),\ ((i,8),(i,9)):0\le i\le8\}.
 \label{rev4:ae:enlaces}
\end{align}
The coordinates are interpreted modulo twenty-seven. The two sets
are disjoint and contain eighteen links each: two sides of
nine links per direction. The labels \(90\) and \(120\) arise
from the incidence spaces; the cardinalities of the link
sets are \(18\) and \(18\).

Let \(u_m\) be the normalized uniform vectors of
\(\ell^2(\mathcal F_6)\oplus\ell^2(\mathcal F_8)\), with
\(m=90,120\), and let \(v_m\) be the corresponding normalized
uniform vectors of \(\ell^2(E_{90})\oplus\ell^2(E_{120})\).
For example, \(u_{90}=90^{-1/2}\sum_{f\in\mathcal F_6}(\delta_f,0)\)
and \(v_{90}=18^{-1/2}\sum_{e\in E_{90}}(\delta_e,0)\).
Each pair is orthonormal.

\begin{proposition}[Preservation of the collective spectral structure]
\label{rev4:ae:colectivo}
The linear map \(\mathcal J u_m=v_m\) is a surjective
isometry between the two collective planes. If
\[
 D_{\rm inc}=90|u_{90}\rangle\langle u_{90}|
             +120|u_{120}\rangle\langle u_{120}|,
 \qquad
 D_\partial=90|v_{90}\rangle\langle v_{90}|
             +120|v_{120}\rangle\langle v_{120}|,
\]
then \(\mathcal J D_{\rm inc}=D_\partial\mathcal J\).
The sectoral projectors and their oriented combinations are
transported by the same map.
\end{proposition}
\begin{proof}
The map takes one orthonormal basis to another and is determined
by those images. The operator equality holds on \(u_m\),
since both sides have value \(m v_m\). The projectors are
\((120I-D)/30\) and \((D-90I)/30\) in each plane, so that
their combinations are preserved by linear composition.
\end{proof}

The preceding transport acts on the uniform modes. The intrinsic
fluctuations of each incidence space remain in their orthogonal
complements. Area weighting uses precisely the two labels
preserved by this isometry.

The Barbero--Immirzi parameter is determined from those data and
the previously constructed angular orientation. Let \(A\) and \(C^*\)
be the angular coordinates in degrees, in the lift fixed in the
preceding section, and set
\[
 x=\frac{\pi_{\rm HMT}A}{180},\qquad
 y=\frac{\pi_{\rm HMT}C^*}{180},\qquad
 q_\pm=e^{-x\mp y},\qquad
 S_m(q)=\frac{q^m}{1-q^{3m}}.
\]
The chamber \(x>|y|\) is preserved, so that \(0<q_\pm<1\).
The oriented chain of one turn assigns twelve sectors to
\(\mathcal F_6\) and an oppositely oriented return seam
to \(\mathcal F_8\). Its evaluation is \(12S_{90}-S_{120}\).
The coefficient of \((1-q)^{-1}\) is determined by
\[
 \lim_{q\uparrow1}(1-q)(12S_{90}(q)-S_{120}(q))
 =\frac{12}{270}-\frac1{360}=\frac1{24},
\]
because \(1-q^{3m}=(1-q)\sum_{j=0}^{3m-1}q^j\).
The conjugate evaluation of the chain, together with the angular component,
defines the structural parameter
\begin{equation}
 \gamma_{\rm HMT}
 =\frac{\pi_{\rm HMT}AC^*}{180}
 +12\bigl[S_{90}(q_-)-S_{90}(q_+)\bigr]
 -\bigl[S_{120}(q_-)-S_{120}(q_+)\bigr].
 \label{rev4:ae:gamma}
\end{equation}
The origin of the coefficients \(12:-1\) lies in the oriented chain.
The preceding factorization determines its pole coefficient.
The inversion \(C^*\mapsto-C^*\) interchanges the channels and changes
the sign of \(\gamma_{\rm HMT}\). In what follows, we retain the
branch \(C^*>0\), \(\gamma_{\rm HMT}>0\), with the area
normalization fixed by that same orientation.

\subsection{Longitudinal return and determination of the area scale}
\label{rev5:longitud}

The dimensional unit of area admits an explicit construction
from the same enriched state. We write \(L_\alpha\) for the
resulting physical length; the integer \(L\) denoting the length
of a record in Section~\ref{sec:compat-lectores-formas} retains
its combinatorial meaning. This distinction separates the cardinality
of a representation, the extent of an edge, and the radius of its realisation.

The composition uses the regional coordinates and the record already
generated by APP--TRIT--TPK. APP retains both sheets with their remainders
and quotients; TRIT orients the local regime; TPK prolongs the state,
retaining carry, boundary, route, and memory. The chain
\(w_6\to w_{12}\to w_{18}\to w_{24}\to w_{30}\to R_{36}\to G_9\)
and the order \(U_t\to\Gamma_9\to K_{\rm ph}\) retain
phase return with memory advancement. The joint construction
of the continuum retains its spectral, solenoidal, gauge,
cohomological, and determinantal aspects. The length constructed below
is a subsequent dimensional reading of that structure and its correlated outputs
\(\pi_{\rm HMT},\varphi_{\rm HMT},e_{\rm HMT},\alpha_{\rm HMT}\).

Fix the marked record
\[
 \Omega=\{(j,k,q,u,v):j,k\in\mathbb Z/12\mathbb Z,\quad
 k-j\in\{0,3,6,9\},\quad1\le q\le8,\quad
 u<v,\quad u,v\in\{1,2,4,5,7,8\}\}.
\]
Subtraction is taken in that cyclic group. Its cardinality is
\(N_\Omega=12\cdot4\cdot8\binom62=5760\).
The factor \(120\) comes from the previously constructed octadic
point--pair incidence space; the twelve positions and the four
positions in each orbit retain their respective markings. This
resolution extends the carrier of the Hadamard return. Between the
initial and final Hermitian spaces, each of dimension \(N_\Omega\),
let \(B\) be the transported inverse, with
\(B^*B=BB^*=I/4\). The local identity
\(H_4^*H_4=4I\) produces this normalisation for the inverse;
its orthogonal extension and the transport of the markings preserve it.

In the final space, let
\(u_\Omega=N_\Omega^{-1/2}\mathbf1\),
\(P=u_\Omega u_\Omega^*\), \(C=(I-P)B\), and \(M=PB\).
The complete transport \(C+M=B\) is retained. The atomic projectors
\(P_i=e_ie_i^*\), determined by the inscription,
define the diagonal expectation
\(\Delta_\Omega(T)=\sum_iP_iTP_i\).

\begin{proposition}[Recorded response and retained complement]
\label{rev5:retorno}
Among the linear maps into the diagonal algebra that fix
the \(P_i\) and are bimodular over that algebra,
\(\Delta_\Omega\) is unique. Moreover,
\begin{equation}
 \Delta_\Omega(CC^*)=r_\Omega I,
 \qquad r_\Omega=\frac{N_\Omega-1}{4N_\Omega}
 =\frac{5759}{23040},
 \qquad \Delta_\Omega(MM^*)=\frac{I}{4N_\Omega}.
 \label{rev5:retorno-escalar}
\end{equation}
The two responses sum to \(I/4\).
\end{proposition}
\begin{proof}
For the matrix unit \(E_{ij}\), bimodularity imposes
\(E(E_{ij})=P_iE(E_{ij})P_j\). Since the image is diagonal,
this term is zero when \(i\ne j\); when \(i=j\), fixing
\(P_i\) determines its image. The matrix units form a basis
and determine the map. Furthermore,
\(CC^*=(I-P)/4\), \(MM^*=P/4\), and every diagonal entry
of \(P\) is \(1/N_\Omega\). Applying the expectation proves
the three identities.
\end{proof}

The isometric inscription \(We_i=e_i\otimes e_i\) realises
\(W^*(T\otimes I)W=\Delta_\Omega(T)\). The complete inscription
retains the information in the record; the expectation extracts
its diagonal response. The complement \(M\) remains in the
transport and has the specific response given in
\eqref{rev5:retorno-escalar}. This separation identifies the observable
entering the length and the complementary component that is retained.

The Smith normal form of \(H_4\) is
\(\operatorname{diag}(1,2,2,4)\); its cokernel has order
sixteen. The product of the generated scalar
\(\alpha=\alpha_{\rm HMT}\), indexed by the classes of that
cokernel, defines the finite multiplicative reader
\[
 n_H:=\prod_{\bar v\in\operatorname{coker}_{\mathbb Z}(H_4)}
 \alpha=\alpha^{16}.
\]
Let \(L_*>0\) be a radial reference
section of the dimensional line of length. The longitudinal composition
acts linearly on an oriented cochain \(\beta_*\):
\begin{equation}
 \mathscr T_L\beta_*
 =\alpha^{16}(\Delta_\Omega(CC^*)\otimes I)\beta_*
 =q_\alpha\beta_*,\qquad
 q_\alpha=\alpha^{16}r_\Omega,\qquad
 L_\alpha=q_\alpha L_*.
 \label{rev5:longitud-generada}
\end{equation}
The formula specifies a concrete composition of readers. Linearity
preserves the length type and its orientation; the coefficient comes
from the recorded return and the local norm. For example, the diagonal
response to a source \(a e_i\) is \(a r_\Omega\), whereas
the norm \(\|C^*ae_i\|=|a|\sqrt{r_\Omega}\) constitutes
another reading of the same transport.

If \(\beta_*(\bar a)=-U(a)\beta_*(a)\), multiplication by
\(q_\alpha\) preserves oriented inversion. If an edge is
subdivided into nine compatible segments,
\(\beta_*(a)=\sum_jU_j^{-1}\beta_*(a_j)\), the same
multiplication preserves that identity and memory concatenation.
In an auxiliary extension of the record,
\(P\), \(CC^*\), and \(\Delta_\Omega\) are transported
by tensoring with the identity. Thus \(r_\Omega\) is preserved;
the cardinality of a refined mesh and the cardinality of the marked
record have distinct functions.

\subsection{Radial reciprocity and recognition of the Planck length}
\label{rev5:radial}

The circular reading of the one-hundred-and-eight-step calendar retains
\(108\ell_0=2\pi_{\rm HMT}L_\alpha\), with
\(\ell_0=ct_0\). Consequently,
\begin{equation}
 t_0=\frac{\pi_{\rm HMT}L_\alpha}{54c},\qquad
 \omega=\frac{c}{L_\alpha},\qquad
 Q_{\rm clk}=\hbar_{\rm ret}\omega
 =\frac{\hbar_{\rm ret}c}{L_\alpha}.
 \label{rev5:reloj}
\end{equation}
The returned action is obtained from its full HMT equation in
\eqref{eq:mass-k-action}; the speed is obtained from the two constitutive
channels of Section~\ref{sec:compat-lectores-formas}.
The circular relation converts the lifted advance into a radius. If
\(c=\widehat c\,u_C\), \(t_0=u_T\), and
\(u_L=u_Cu_T\), then
\(L_*/u_L=54\widehat c/(\pi_{\rm HMT}q_\alpha)\).
The radial reference \(L_*\) and the edge unit \(u_L\)
are related with their dimensional types made explicit.

In a finite fibre, let \(X=X^*>0\) be an invertible positive
character, \(U=2B\), and \(Y=UXU^*\). The longitudinal
transport is \(D_L=L_\alpha U\). Its positive radial response
is defined by the metric in the final space,
\(\|R_Xv\|^2=\|XD_L^*v\|^2\) for every \(v\).
The energy and conjugate length retain the same character:
\[
 H_X=Q_{\rm clk}Y,\qquad M_X=H_X/c^2,\qquad
 \Lambda_X=L_\alpha Y^{-1}.
\]

\begin{theorem}[Common radial coefficient and area scale]
\label{rev5:rigidez-radial}
The preceding conditions determine a unique positive response
\(R_X=L_\alpha Y\). The following identities hold:
\begin{equation}
 R_X\Lambda_X=L_\alpha^2I,\qquad
 H_X\Lambda_X=\hbar_{\rm ret}cI,\qquad
 R_X=\frac{L_\alpha^2}{\hbar_{\rm ret}c}H_X.
\end{equation}
In the physical realisation identifying \(R_X\) with the reduced
gravitational radius \(R_X=(G/c^2)M_X\), the coefficient is
unique and common to all admissible positive characters:
\begin{equation}
 G=\frac{c^3L_\alpha^2}{\hbar_{\rm ret}},\qquad
 \frac{\hbar_{\rm ret}G}{c^3}=L_\alpha^2.
 \label{rev5:G-area}
\end{equation}
\end{theorem}
\begin{proof}
Polarisation of the equality of norms determines
\(R_X^2=D_LX^2D_L^*=L_\alpha^2UX^2U^*\).
The operator \(L_\alpha UXU^*\) is positive and its square
is the specified operator. Uniqueness of the positive square root
proves the first assertion. The two product identities follow by
substitution. The physical convention gives
\(L_\alpha Y=(GQ_{\rm clk}/c^4)Y\). Invertibility of
\(Y\) allows its cancellation, and the expression for
\(Q_{\rm clk}\) yields \eqref{rev5:G-area}.
If another coefficient preserves the same data, its difference
multiplied by \(M_X\) is zero and is therefore itself zero.
\end{proof}

The proof applies to the finite positive fibre; it also extends to
bounded uniformly positive operators satisfying the same identities.
The reduced-gravitational-radius reading is the declared physical
convention of this realisation. The subsequent conventional recognition
\(\ell_P^2=\hbar_{\rm ret}G/c^3\) identifies
\(\ell_P=L_\alpha\). The length is first obtained through
\eqref{rev5:longitud-generada}; this identification preserves its
construction. The equivalent circular formula is
\[
 G=\left(\frac{54}{\pi_{\rm HMT}}\right)^2
 \frac{c^5t_0^2}{\hbar_{\rm ret}}.
\]
The circular factor is already included when \(L_\alpha\) is used.
The length and action section also determine
\[
 t_P=\frac{L_\alpha}{c},\qquad
 m_P=\frac{\hbar_{\rm ret}}{cL_\alpha},\qquad
 E_P=\frac{\hbar_{\rm ret}c}{L_\alpha}.
\]
These identities follow by substituting \eqref{rev5:G-area} into
the conventional expressions for the three scales; positivity
determines their roots. For a mode of character \(y>0\), the
readings are \(m(y)=m_Py\), \(r_g(y)=L_\alpha y\), and
\(\bar\lambda(y)=L_\alpha/y\). Their longitudinal product
is \(r_g(y)\bar\lambda(y)=L_\alpha^2\). Equality of
the two lengths is equivalent to \(y=y^{-1}\) and, by positivity,
characterises exactly \(y=1\). The calendar retains
\(t_0=(\pi_{\rm HMT}/54)t_P\) and
\(108t_0=2\pi_{\rm HMT}t_P\).

A horizon radius \(R_h=L_\alpha\zeta_h\) retains its own
factor \(\zeta_h=R_h/L_\alpha\). Its spatial area belongs
to the specified horizon; the elementary scale of the following
sectorial operator is \(L_\alpha^2\). In the Schwarzschild
specialisation for the same mass, the radius is \(2r_g(y)\)
and thus \(\zeta_h=2y\); a horizon in another realisation
retains its own geometric formation condition.


\subsection{Boundary occupations and sectoral reconstruction}

Each link of \(E_{90}\sqcup E_{120}\) determines a factor
\(\mathbb C^3\), with basis \(|0\rangle,|1\rangle,|2\rangle\)
and number operator \(N_e=\operatorname{diag}(0,1,2)\). The boundary
state space is
\(\mathscr H_A=(\mathbb C^3)^{\otimes18}\otimes
(\mathbb C^3)^{\otimes18}\). We define \(N_m\) as the sum of
the number operators of the eighteen links in sector \(m\).
The two operators commute and each has integer spectrum between
zero and thirty-six.

The angular scale arises from the clock and the normalized character of the
conjugate triplet:
\begin{equation}
 \theta_{\rm clk}=\frac{C^*}{6}\frac{\pi_{\rm HMT}}{180},
 \quad
 \chi_{10}=\frac13\Tr\operatorname{diag}
 (1,e^{i\pi_{\rm HMT}/18},e^{-i\pi_{\rm HMT}/18})
 =\frac{1+2\cos(\pi_{\rm HMT}/18)}3,
 \quad a_0=\chi_{10}\theta_{\rm clk}.
 \label{rev4:ae:escala}
\end{equation}
The character identity is obtained by summing the conjugate phases.
On the branch considered, \(a_0>0\). The area unit is
\(L_\alpha^2=\ell_P^2\), constructed and recognized in
\eqref{rev5:longitud-generada}--\eqref{rev5:G-area}. Write
\(a_\partial=\gamma_{\rm HMT}a_0>0\). The area operator is
\begin{equation}
 \widehat{\mathcal A}=L_\alpha^2a_\partial N,
 \qquad N=N_{90}+N_{120},
 \qquad D=90N_{90}+120N_{120}.
 \label{rev4:ae:area}
\end{equation}
Its spectrum is \(\{nL_\alpha^2a_\partial:0\le n\le72\}\).
The multiplicity of \(n\) is the coefficient of \(z^n\) in
\((1+z+z^2)^{36}\), since each tensor factor contributes one of
the three occupations.

\begin{theorem}[Reconstruction of the two sectoral occupations]
\label{rev4:ae:recuperacion}
The pair of observables \((N,D)\) determines
\begin{equation}
 N_{90}=4N-\frac{D}{30},
 \qquad N_{120}=\frac{D}{30}-3N.
 \label{rev4:ae:inversa}
\end{equation}
In the joint spectrum, their values \((n,d)\) satisfy
\(n\in\mathbb Z\), \(d\in30\mathbb Z\), \(90n\le d\le120n\) and
\(0\le4n-d/30\le36\), \(0\le d/30-3n\le36\).
These conditions characterize the image of the integer occupations.
\end{theorem}
\begin{proof}
The transformation from \((N_{90},N_{120})\) to \((N,D)\) has matrix
\(\left(\begin{smallmatrix}1&1\\90&120\end{smallmatrix}\right)\),
whose determinant is \(30\). Its inverse gives the operator
identities. Commutativity allows them to be applied in each joint
eigenspace. The inequalities express nonnegativity and the maximum
capacities of the sectors; divisibility makes both occupations integers.
Conversely, each integer between zero and thirty-six can be expressed as
a sum of eighteen terms from \(\{0,1,2\}\), and therefore every pair
satisfying these conditions belongs to the joint spectrum.
\end{proof}

Area preserves the sum of occupations; the weighted observable preserves
their sectoral composition. Thus, \((N_{90},N_{120})=(1,0)\) and \((0,1)\)
have the same area \(L_\alpha^2a_\partial\) and weights \(90\) and \(120\).
Complementarily, \((4,0)\) and \((0,3)\) have the same weight
\(D=360\), but different areas. The pair distinguishes both situations.
Within each joint eigenspace, the microscopic degeneracies
corresponding to the distributions among links are preserved.

\begin{figure}[htbp]
\centering
\begin{tikzpicture}[>=Latex,every node/.style={font=\small},node distance=15mm]
\node[draw,rounded corners,align=center] (a) {Sectoral occupation\\$(1,0)$};
\node[draw,rounded corners,align=center,right=54mm of a] (b) {Sectoral occupation\\$(0,1)$};
\node[draw,rounded corners,align=center] (n) at ($(a)!0.5!(b)+(0,-1.8)$)
 {Common area\\$L_\alpha^2a_\partial$};
\draw[->] (a)--node[left] {$N=1$}(n);
\draw[->] (b)--node[right] {$N=1$}(n);
\node[above=5mm of a] {$D=90$};
\node[above=5mm of b] {$D=120$};
\end{tikzpicture}
\caption{Equality of area preserves total occupation. Incidence
weighting distinguishes the two compositions; their joint reading allows
the occupations of both sectors to be reconstructed.}
\label{rev4:ae:figura}
\end{figure}

\subsection{Reduced state, purification and entanglement entropy}

The informational interpretation requires the state to be fixed. For
finite \(\Delta\in\mathbb R\), with the normalization of the modular
parameter of the sectoral construction, we define
\begin{equation}
 x_m=e^{-m\Delta},\qquad z_m=1+x_m+x_m^2,
 \qquad Z=z_{90}^{18}z_{120}^{18},
 \qquad \rho_A(\Delta)=Z^{-1}e^{-\Delta D}.
 \label{rev4:ae:rho}
\end{equation}
The sum of commuting occupations and the tensor decomposition prove
the expression for \(Z\). All eigenvalues of \(\rho_A\) are
strictly positive for finite real \(\Delta\). Its modular
Hamiltonian, denoted by \(K_A\) to distinguish it from the arithmetic record
\(K\), is
\begin{equation}
 K_A=-\log\rho_A=(\log Z)I+\Delta D.
 \label{rev4:ae:modular}
\end{equation}
This operator characterizes the spectral distribution of the reduced state.
The time-evolution generator is defined in the dynamical development
with its own parameters and domain. For \(\Delta\ne0\), with
\(\Delta\) and the normalization known, the pair
\((\widehat{\mathcal A},K_A)\) recovers \((N,D)\) and, through
\eqref{rev4:ae:inversa}, both occupations. For \(\Delta=0\),
the state is uniform and \(K_A=36\log3\,I\).

We construct the bipartition explicitly. Let \(\mathscr H_{\bar A}\)
be a copy of \(\mathscr H_A\), with the conjugate bases fixed. For
a link of weight \(w\), the vector
\[
 |\operatorname{TFD}(w)\rangle
 =\frac1{\sqrt{1+e^{-w}+e^{-2w}}}
 \sum_{n=0}^{2}e^{-wn/2}|n\rangle_A\otimes|n\rangle_{\bar A}
\]
has norm one. The pure state
\begin{equation}
 |\Psi_\Delta\rangle
 =|\operatorname{TFD}(90\Delta)\rangle^{\otimes18}
  \otimes|\operatorname{TFD}(120\Delta)\rangle^{\otimes18}
 \label{rev4:ae:tfd}
\end{equation}
determines the bipartition \(A\mid\bar A\).

\begin{proposition}[Sectoral reduction and entropy of the purification]
\label{rev4:ae:entropia}
The partial trace of \(|\Psi_\Delta\rangle\langle\Psi_\Delta|\)
over \(\mathscr H_{\bar A}\) is \(\rho_A(\Delta)\). Defining
\[
 f_m=\frac{x_m+2x_m^2}{z_m},
 \qquad s(\Delta)=-\Tr(\rho_A(\Delta)\log\rho_A(\Delta)),
\]
we have
\begin{align}
 \langle N\rangle_\Delta&=18(f_{90}+f_{120}),\\
 \langle D\rangle_\Delta&=18(90f_{90}+120f_{120}),\\
 \langle\widehat{\mathcal A}\rangle_\Delta
 &=18L_\alpha^2a_\partial(f_{90}+f_{120}),\\
 s(\Delta)&=18\log z_{90}+18\log z_{120}
             +\Delta\langle D\rangle_\Delta.
 \label{rev4:ae:entropia-formula}
\end{align}
The last expression is the entanglement entropy of the
purification \eqref{rev4:ae:tfd}, in dimensionless units.
\end{proposition}
\begin{proof}
In the partial trace of a link, the orthogonality of
\(|n\rangle_{\bar A}\) eliminates terms with distinct indices and
leaves the weights \((1,x_m,x_m^2)/z_m\). The partial trace of a tensor
product is the product of the partial traces; this yields
\eqref{rev4:ae:rho}. The first moment of each link is \(f_m\).
Additivity of the number operators gives the first three identities.
The last follows by substituting \eqref{rev4:ae:modular} into
\(s=\Tr(\rho_AK_A)\). The coefficients of the preceding construction
are Schmidt coefficients of the stated bipartition, so
this entropy coincides with its entanglement entropy.
\end{proof}

The maximum entropy of this space is \(36\log3\), attained at
\(\Delta=0\). The quantity \(36\log3-s(\Delta)\) is the
relative entropy of \(\rho_A(\Delta)\) with respect to
\(I/3^{36}\). Incidence fixes the sectors and their multiplicities;
the state parameter fixes the probabilities; the dimensional chart
converts occupation into area. The physical identification of
\(\mathscr H_{\bar A}\) preserves the bipartition and the algebra of
observables just specified.

\subsection{Finite modular identity and transport of geometric information}

Fix the reference state \(\rho_0=\rho_A(\Delta_0)\), with
finite real \(\Delta_0\), and hold
\(\Delta_0,a_\partial,L_\alpha\) constant. The normalized sectoral
areas \(\mathcal A_m^{\rm n}=a_\partial N_m\) satisfy
\begin{equation}
 K_0=-\log\rho_0=c_0I+
 \beta_{90}\mathcal A_{90}^{\rm n}
 +\beta_{120}\mathcal A_{120}^{\rm n},
 \qquad c_0=\log Z(\Delta_0),\qquad
 \beta_m=\frac{m\Delta_0}{a_\partial}.
 \label{rev4:ae:betas}
\end{equation}
The relation \(\beta_{120}=\tfrac43\beta_{90}\) preserves the incidence
ratio. For \(\Delta_0\ne0\), the ratio of the two coefficients
is defined and equals \(4/3\).

\begin{theorem}[Finite identity of area and relative information]
\label{rev4:ae:ley-finita}
For every normalized state \(\sigma\) of \(\mathscr H_A\), let
\(\delta_\sigma\langle B\rangle=\Tr[(\sigma-\rho_0)B]\).
Then
\begin{equation}
 s(\sigma)-s(\rho_0)
 =\sum_{m\in\{90,120\}}\beta_m
       \delta_\sigma\langle\mathcal A_m^{\rm n}\rangle
   -D(\sigma\Vert\rho_0),
 \label{rev4:ae:ley-finita-formula}
\end{equation}
where \(D(\sigma\Vert\rho_0)=
\Tr[\sigma(\log\sigma-\log\rho_0)]\) is finite and nonnegative,
with the convention \(0\log0=0\).
\end{theorem}
\begin{proof}
The definition yields
\(D(\sigma\Vert\rho_0)=-s(\sigma)+\Tr(\sigma K_0)\),
while \(s(\rho_0)=\Tr(\rho_0K_0)\). Subtracting the identities
and substituting \eqref{rev4:ae:betas} proves
\eqref{rev4:ae:ley-finita-formula}; the scalar term vanishes
by normalization of the traces. For positivity, let \(p_i\)
and \(u_i\) be the eigenvalues and eigenvectors of \(\sigma\), and
\(q_j>0,v_j\) those of \(\rho_0\). With
\(r_i=\langle u_i,\rho_0u_i\rangle\), concavity of
\(\log\) gives
\(\langle u_i,\log\rho_0u_i\rangle\le\log r_i\).
Thus, \(D(\sigma\Vert\rho_0)\ge\sum_i p_i\log(p_i/r_i)\ge0\),
by convexity of \(t\log t\), since \(\sum_i r_i=1\).
Full rank of \(\rho_0\) guarantees finiteness.
\end{proof}

Its infinitesimal version keeps the reference state fixed. If
\(\sigma(\varepsilon)\) is differentiable, normalized and
\(\sigma(0)=\rho_0\), then
\[
 \left.\frac{\dd s(\sigma(\varepsilon))}{\dd\varepsilon}\right|_0
 =\sum_m\beta_m
  \left.\frac{\dd\langle\mathcal A_m^{\rm n}\rangle_{
                          \sigma(\varepsilon)}}{\dd\varepsilon}\right|_0.
\]
Indeed, the derivative of the trace of \(-\sigma\log\sigma\) is
\(-\Tr[\dot\sigma(\log\rho_0+I)]\); normalization eliminates
\(\Tr\dot\sigma\) and substitution of \(K_0\) gives the formula.
The finite identity adds the exact relative-entropy correction
to that tangent response. The coefficients of the physical areas are
\(\beta_m/L_\alpha^2\), preserving the same dimensional chart.

Composition with the Grassmannian representation finally specifies
what information is preserved. In
\(\mathcal F(K,\xi)=([B(K)],Q,\xi)\), the additional coordinates
\(\xi\) retain the regional magnitudes, angular orientation,
marked inclusion, sector identification, occupation
basis, parameter \(\Delta\) and relevant area unit.
The inverse of Proposition~\ref{prop:compat-inversa-marcada}
recovers \(K\); the other conditions are preserved through their
markings. By Theorem~\ref{thm:compat-algebra-lectores}, composition
with that inverse transports \(\gamma_{\rm HMT}\),
\(\widehat{\mathcal A}\), \(\rho_A\), \(K_A\) and their evaluations
on the stated domain. Changes of chart of the same marked point
preserve these objects, and refinements preserve the inherited
evaluations when they maintain the same sectoral data.

This articulation preserves three distinguishable levels of information.
The marked positive representation recovers the record with the retained
data; the area--modular-Hamiltonian pair recovers the two
occupations under the conditions of \eqref{rev4:ae:modular}; the
purification determines the entanglement of a specific bipartition.
Information loss in a projection is computed through its fibers:
area identifies occupations with equal sums, while the joint
reading distinguishes their composition and preserves the remaining microscopic
degeneracy. The link between geometry and information is thus expressed
by verifiable maps and identities on the same construction.
\subsection{Naturality of the scale and operator-valued canonical forms}
\label{rev5:naturalidad-areal}

The generated length makes the compatibility between the positive
realisation, the area operator, and sectorial information explicit.
The marked representation \(\mathcal F(K,\xi)\) retains,
among the relevant data in \(\xi\), the resolution
\(\Omega\), its projectors, the transported return, and the section
\(L_*\). The generated regional coordinates, the action, and the constitutive
channels retain their preceding genealogy. The marked inverse
therefore recovers every argument of
\eqref{rev5:longitud-generada} and \eqref{rev5:G-area}.
In particular, the previously constructed geometric reader of
\(\alpha^{\mathrm{geom}}\) determines the expressions
\[
 L_\alpha^{\mathrm{geom}}
 =L_*(\alpha^{\mathrm{geom}})^{16}\frac{5759}{23040},
 \qquad
 G^{\mathrm{geom}}
 =\frac{(c^{\mathrm{geom}})^3
 (L_\alpha^{\mathrm{geom}})^2}{\hbar_{\rm ret}^{\mathrm{geom}}}.
\]
The speed and action readers are obtained by the same composition
with the marked inverse. These expressions retain the return data
and the radial reference section.

\begin{corollary}[Preservation under changes of chart and inherited refinement]
\label{rev5:conservacion}
Positive changes of chart preserving the marked point, and
refinements recovering the initial record, preserve
\(L_\alpha\), \(G\), \(\gamma_{\rm HMT}\), and
\(\widehat{\mathcal A}=L_\alpha^2a_\partial N\), provided
they also transport the declared sectorial and dimensional data.
The normalised spectrum
\(\widehat{\mathcal A}/L_\alpha^2=a_\partial N\), its
multiplicities, and the state \(\rho_A(\Delta)\) remain unchanged.
\end{corollary}
\begin{proof}
Theorem~\ref{thm:compat-algebra-lectores} transports each reader
by composition with the inverse. Mutations of the same point
reconstruct the same record. Under inherited refinement, summing
the subdivided separations reconstructs each initial coordinate.
Tensor extension of the inscription preserves \(r_\Omega\),
while its remaining arguments are fixed by the markings.
Substitution therefore preserves the length, the gravitational
coefficient, and the area operator. The occupation projectors
preserve the multiplicity polynomial \((1+z+z^2)^{36}\) and
the normalised expression of the state.
\end{proof}

The inherited refinement considered here retains the same space
of thirty-six boundary links. An extension introducing a different
physical boundary space requires specified intertwining maps for
\(N_{90}\) and \(N_{120}\), together with the transport of the
declared sectorial data.

Let \(P\) now be a rank-one polytope constructed in this manuscript,
with an oriented subdivision \(P=\bigcup_\sigma P_\sigma\)
having the same support, multiplicity one, and compatible faces.
The same sectorial fibre \(\mathscr H_A\) is fixed over all
its cells. The operator-valued top-degree form with the
dimension of area is defined by
\[
 \boldsymbol\Omega^{\mathcal A}_P
 =\widehat{\mathcal A}\otimes\Omega_P.
\]
Here \(\widehat{\mathcal A}\) is constant with respect to the
differential coordinates of the polytope: it is determined by the
fixed HMT state, whereas \(\Omega_P\) describes its geometric support.
The coefficient \(\widehat{\mathcal A}\) is positive semidefinite.
Its component in the vacuum sector \(N=0\) vanishes. On an
eigenspace with occupation \(n>0\), division by
\(nL_\alpha^2a_\partial\) recovers the canonical form
\(\Omega_P\) multiplied by the identity on that eigenspace.

\begin{proposition}[Areal additivity and boundary residues]
\label{rev5:forma-areal}
In the common sectorial trivialisation,
\begin{equation}
 \boldsymbol\Omega^{\mathcal A}_P
 =\sum_\sigma\widehat{\mathcal A}\otimes\Omega_{P_\sigma},
 \qquad
 \operatorname{Res}_F\boldsymbol\Omega^{\mathcal A}_P
 =\widehat{\mathcal A}\otimes\Omega_F.
\end{equation}
These identities persist under successive subdivisions and
iterated residues. For a fixed state \(\rho\), taking the trace
gives \(\Tr(\rho\widehat{\mathcal A})\Omega_P\) and preserves
the same scalar identities.
\end{proposition}
\begin{proof}
The operator space of \(\mathscr H_A\) is finite-dimensional.
In any basis, the assertion is the additivity identity for
canonical forms multiplied by each constant matrix entry of
\(\widehat{\mathcal A}\). Both residue and trace are linear.
Cancellation at interior facets retains their orientations and
the same operator entry on both sides. Iteration applies these
operations to each subdivision and each flag of faces.
\end{proof}

If distinct cells carry regular operator-valued coefficients
\(A_\sigma\), the defect of a pair along its common face is
\((A_\sigma|_F-A_\tau|_F)\otimes\Omega_F\).
Equality of the operator-valued restrictions to the boundary
provides the gluing condition, under the simple-pole hypotheses
of Theorem~\ref{thm:compat-pegado}. The global sum must also retain
every contribution along each ambient divisor and the control
of exterior poles. Dimensional normalisation, residue cancellation,
and coverage of the support therefore have distinct and
compatible criteria.

The finite modular identity likewise admits a direct expression
in physical areas \(\mathcal A_m=L_\alpha^2a_\partial N_m\):
\begin{equation}
 s(\sigma)-s(\rho_0)
 =\sum_{m=90,120}\frac{m\Delta_0}{L_\alpha^2a_\partial}
   \Tr[(\sigma-\rho_0)\mathcal A_m]
  -D(\sigma\Vert\rho_0).
 \label{rev5:modular-fisica}
\end{equation}
This is \eqref{rev4:ae:ley-finita-formula} after substitution
of the constructed length. The angular factor \(a_0\), contained
in \(a_\partial=\gamma_{\rm HMT}a_0\), is retained in full.
When dimensional sections satisfying \(L'_\alpha=sL_\alpha\)
are compared while preserving the same state and the same
dimensionless data, areas are multiplied by \(s^2\) and their
modular coefficients by \(s^{-2}\). Entropy and relative entropy
remain unchanged. This comparison of sections does not constitute
a second solution for the same fixed longitudinal data.

For a spherical horizon of radius \(R_h=L_\alpha\zeta_h\),
the geometric area is \(4\pi_{\rm HMT}L_\alpha^2\zeta_h^2\).
Its identification with the expectation of the sectorial operator,
when required for the same cross-section, imposes the explicit
condition
\[
 a_\partial\Tr(\rho N)=4\pi_{\rm HMT}\zeta_h^2.
\]
The horizon geometry and the boundary state must satisfy this
relation through their respective constructions. The factor
\(\zeta_h\) characterises that horizon; the longitudinal scale
of the coupling remains \(L_\alpha\).

